We introduced a general diffusion-based framework for stochastic optimisation
when the loss gradient is an expectation under a parameter-dependent
distribution. The framework couples an optimiser with a continuously evolving
sampler and accommodates a wide range of dynamics within a common mean-field formulation. Under
well-posedness, existence of the stationary pair, and twisted joint strong
monotonicity, the mean-field system contracts exponentially to that stationary
solution. For the interacting particle system, we
obtained an explicit \(N^{-1/2}\) interaction-error bound and a parameter convergence bound.

The experiments illustrate the practical flexibility of the framework. In
Bayesian image deblurring, the higher-order Langevin variant improves the
reported MSE and SSIM and produces sharper reconstructions than the proximal
overdamped baseline. In the two-dimensional energy-based-model
experiments, the momentum method reaches lower Fisher divergence more rapidly
under the matched (gradient-evaluation) budget, while retaining
comparable visual sample quality.

\subsection{Limitations and Future Work}

The theory relies on global contractivity and uniform-moment assumptions, which
can be restrictive for the non-convex objectives encountered in practice. Moreover, the guarantees concern the continuous-time systems:
establishing finite-step error and stability for the proposed splitting schemes
requires a separate discretisation analysis. The time-dependent contraction
result is conditional on verifying its monotonicity rate for each schedule. Future work should therefore develop local or non-convex guarantees, analyse discretisation and adaptive preconditioning.